% Technical Report: Network-Topology Correction for Accessibility Modelling
% Mirrors the preamble of paper/main.tex (LLNCS format).
%
\documentclass{llncs}

% ─── Encoding & Fonts ─────────────────────────────────────────────────────────
\usepackage[utf8]{inputenc}
\usepackage{lmodern}

% ─── Mathematics ──────────────────────────────────────────────────────────────
\usepackage{amsmath}
\usepackage{amssymb}

% ─── Tables ───────────────────────────────────────────────────────────────────
\usepackage{array}
\usepackage{booktabs}
\usepackage{tabularx}

% ─── Algorithms ───────────────────────────────────────────────────────────────
\usepackage{algorithm}
\usepackage{algpseudocode}

% ─── Figures & Floats ─────────────────────────────────────────────────────────
\usepackage{graphicx}
\usepackage{float}

% ─── Layout & Formatting ──────────────────────────────────────────────────────
\usepackage{indentfirst}

% ─── References & Links ───────────────────────────────────────────────────────
\usepackage{cite}
\usepackage{hyperref}

% ─── Custom commands (from main.tex) ──────────────────────────────────────────
\newcommand{\quo}[1]{``#1''}
\newcommand{\mycomment}[1]{}

\begin{document}

\title{Technical Note: Network-Topology-Aware Accessibility Modelling\\
       for the MO-IHLNDP Synthetic Dataset}

\titlerunning{Network-Topology Correction for MO-IHLNDP Dataset}

\author{
    Phu-Quy Le-Tang\inst{1} \and
    Quang-Vu Nguyen\inst{1} \and
    Sunarin Chanta\inst{2}
}

\authorrunning{Le Tang Phu Quy et al.}

\institute{
    The University of Danang -- Vietnam-Korea University of Information and
    Communication Technology
    \and
    King Mongkut's University of Technology North Bangkok \\
    \email{quyltp.22git@vku.udn.vn, nqvu@vku.udn.vn, sunarin.c@itm.kmutnb.ac.th}
}

\maketitle

\begin{abstract}
This technical note documents a methodological correction to the synthetic
dataset used in the MO-IHLNDP case study for Central Vietnam.
The original dataset generated arc-accessibility indicators under a
complete-graph assumption, which conflates link existence with link disruption.
We replace this with a two-stage model: a Delaunay triangulation defines the
planar network topology, and BFS path-reachability extends accessibility to
non-adjacent node pairs.
Transport costs for non-adjacent pairs are updated to shortest-path distances
through the planar network.
This note describes the original methodology, its limitations, the corrective
model, and the implementation plan.
\end{abstract}

\keywords{Synthetic dataset \and Accessibility modelling \and Delaunay
triangulation \and Path reachability \and Humanitarian logistics}


% ============================================================================
\section{Network-Topology Limitation and Planned Correction}
\label{sec:topology-correction}
% ============================================================================

\subsection{Current Methodology}
\label{sec:current}

The synthetic dataset generates the arc-accessibility indicator $a_{uvms}$
under a \textbf{complete-graph assumption}: for every ordered pair
$(u,v)\in V\times V$, across all modes $m$ and scenarios $s$, a disruption
probability is computed as
%
\begin{equation}
  p^{\text{road}}_{uvs}
    = \min\!\bigl(0.97,\;\beta_s \cdot \bar{r}_{uvs}\bigr),
  \qquad
  \bar{r}_{uvs} = \tfrac{r_{us}+r_{vs}}{2},
\label{eq:p_road}
\end{equation}
%
where $\beta_s \in \{0.25, 0.55, 0.88\}$ scales with scenario severity
(mild, severe, extreme).
Water mode is enabled for a pair $(u,v)$ when both endpoints exceed a
flood-risk threshold ($r_{us} > 0.30$ and $r_{vs} > 0.30$); air remains
universally accessible.
Transport costs and travel times are estimated via Haversine distances
adjusted for mode-specific tortuosity and terrain correction factors,
computed independently for all $N^2$ pairs.

This formulation is aligned with prior two-stage stochastic relief logistics
models~\cite{barzinpour2014,tofighi2016humanitarian} that adopt complete-graph
networks for analytical tractability.

\subsection{Identified Limitation}
\label{sec:limitation}

The complete-graph assumption conflates \textbf{link existence} with
\textbf{link disruption}: it assigns a disruption probability to arcs that
have no corresponding physical road or waterway in Central Vietnam's geography.
Two consequences arise.

\textit{First}, nodes on the convex boundary of the study region receive
artificially high connectivity.
Their pairwise risk average $\bar{r}_{uvs}$ is diluted by distant,
lower-risk neighbours, so \quo{belt} arcs along the convex hull remain
accessible even under extreme flooding.
This contradicts real infrastructure patterns along the narrow coastal
corridor of the Quang Nam -- Da Nang -- Thua Thien-Hue provinces, where
coastal roads are typically the first to be inundated or cut~\cite{nguyen2021flood}.

\textit{Second}, transport costs for non-adjacent node pairs underestimate
true distances: direct Haversine costs ignore mandatory detours through the
inland mountain passes of the Truong Son range, whose terrain factor
$\tau(\ell) \in [1.0, 1.8]$ substantially inflates effective travel distance
for nodes separated by highland terrain.


\subsection{Proposed Correction: Two-Stage Topology-Aware Model}
\label{sec:correction}

We replace the complete-graph with a \textbf{two-stage network-topology-aware
accessibility model}, implemented in \texttt{data\_generate\_cv.py}.

\begin{enumerate}

\item \textbf{Planar base graph.}
A Delaunay triangulation of all node coordinates
defines the set of \emph{directly adjacent} arcs
$\mathcal{E} \subset V\times V$, yielding
$|\mathcal{E}| \approx 3N$ edges (346 edges for $N=132$) ---
a sparse, geographically plausible planar network.
The triangulation is computed once per instance via
\texttt{scipy.spatial.Delaunay}.

\item \textbf{Edge-level disruption.}
The risk-based disruption (Eq.~\ref{eq:p_road}) and water-threshold logic
are applied exclusively to $\mathcal{E}$, producing a scenario-specific
accessible subgraph
$G_s^m = (V,\,\mathcal{E}_s^m)$ for each mode $m$.
Non-adjacent arcs are initialised to zero (blocked) and are never assigned
a disruption draw.

\item \textbf{Path-reachability extension.}
For all non-adjacent pairs $(u,v)$, we set
$a_{uvms} = 1$ if and only if $v$ is reachable from $u$
in $G_s^m$ via BFS; otherwise $a_{uvms}$ remains zero.
This preserves the binary structure of the existing constraint set
(Constraints~10--11 in the original model) while capturing realistic
route detours through intermediate nodes.

\item \textbf{Shortest-path transport costs.}
For road (mode~0) and water (mode~1), the direct Haversine cost is retained
for Delaunay-adjacent pairs; for non-adjacent pairs the cost matrix entry
$C_{uvm}$ is overwritten with the Dijkstra shortest-path cost through
$\mathcal{E}$, so that model decisions correctly reflect the cost of the
least-expensive multi-hop route.
Air (mode~2) retains direct Haversine costs, as helicopters fly straight lines.

\end{enumerate}

This correction ensures that under extreme scenarios, isolated coastal nodes
genuinely lose road access rather than retaining it through spurious direct
arcs, and that water transport's compensatory role along the Thu Bon and
Huong rivers is geographically grounded.

\subsection{Quantitative Impact}
\label{sec:impact}

Table~\ref{tab:access} reports the number of non-Delaunay node pairs that
carry a road-accessibility value of one (i.e., path-reachable) under the
corrected model for CV-Large ($N=132$, 8\,300 non-adjacent pairs).

\begin{table}[h]
\centering
\caption{Path-reachable non-adjacent node pairs (road mode, CV-Large).}
\label{tab:access}
\begin{tabular}{lcc}
\toprule
Scenario & Reachable pairs & Fraction \\
\midrule
Mild    & 6\,794 & 81.9\% \\
Severe  & 6\,794 & 81.9\% \\
Extreme & 6\,016 & 72.5\% \\
\bottomrule
\end{tabular}
\end{table}

Under the original complete-graph model, all $N(N-1)/2 = 8\,646$ non-diagonal
pairs received independent disruption draws; under the corrected model, only
the 346 Delaunay edges receive draws, and path-reachability determines the
rest.
The extreme scenario shows a meaningful reduction of 778 additional
unreachable pairs compared with mild/severe, correctly capturing the greater
network fragmentation caused by large-scale flooding.


\subsection{Implementation Plan}
\label{sec:impl}

Algorithm~\ref{alg:accessibility} summarises the corrected accessibility
generation procedure, replacing lines 483--577 of
\texttt{src/scripts/data\_generate\_cv.py}.

\begin{algorithm}[h]
\caption{Topology-aware accessibility generation (per scenario $s$)}
\label{alg:accessibility}
\begin{algorithmic}[1]
\Require Node coordinates $\text{coords}$, scenario risk vector $\mathbf{r}_s$,
         disruption parameter $\beta_s$
\Ensure  Accessibility tensor $a[m][u][v]$ for $m \in \{0,1,2\}$

\State $\mathcal{E} \leftarrow \textsc{DelaunayEdges}(\text{coords})$
       \Comment{$\approx 3N$ undirected edges}
\State $a[m][u][v] \leftarrow 0$ for all $m,u,v$
       \Comment{default: all blocked}

\ForAll{$(u,v) \in \mathcal{E}$}
  \State $\bar{r} \leftarrow (r_{us} + r_{vs})/2$
  \State $a[0][u][v] \leftarrow a[0][v][u] \leftarrow
         \mathbf{1}[\text{Uniform}(0,1) \geq \min(0.97,\,\beta_s \cdot \bar{r})]$
         \Comment{road}
  \State $a[1][u][v] \leftarrow a[1][v][u] \leftarrow
         \mathbf{1}[r_{us} > 0.30 \wedge r_{vs} > 0.30]$
         \Comment{water}
  \State $a[2][u][v] \leftarrow a[2][v][u] \leftarrow 1$
         \Comment{air}
\EndFor

\ForAll{$m \in \{0,1,2\}$}
  \State Build adjacency list $\text{adj}_m$ from $\{(u,v)\in\mathcal{E} : a[m][u][v]=1\}$
  \ForAll{$\text{src} \in V$}
    \State $\text{visited} \leftarrow \textsc{BFS}(\text{adj}_m, \text{src})$
    \ForAll{$\text{dst} \in \text{visited},\; \text{dst} \neq \text{src}$}
      \State $a[m][\text{src}][\text{dst}] \leftarrow 1$
    \EndFor
  \EndFor
\EndFor
\end{algorithmic}
\end{algorithm}

Transport costs for non-adjacent pairs are updated by running Dijkstra from
every source node through $\mathcal{E}$ (road and water modes only),
replacing the direct Haversine entry with the shortest-path distance.
The visualiser (\texttt{visualizer/dataset\_view.py}) requires no changes:
the Delaunay edge display will now correctly reflect the corrected data
topology once the instance JSON files are regenerated.


\bibliographystyle{ieeetr}
\bibliography{cite-class.bib}

\end{document}
